\documentclass[../../../include/open-logic-section]{subfiles}

\begin{document}
\olfileid{sth}{ordinals}{zfm}	
\olsection{$\ZFminus$: یک نقطهٔ عطف}

پرسش از چگونگی توجیه جایگزینی (اگر اصلاً توجیه‌پذیر باشد) سرراست نیست.
ازاین‌رو آن را برای~\olref[replacement][]{chap} کنار می‌گذاریم. اما با
افزودن جایگزینی به نقطهٔ عطف مهم دیگری رسیده‌ایم. اکنون همهٔ اصول موضوع
لازم برای نظریهٔ~$\ZFminus$ را در اختیار داریم. به‌تفصیل:

\begin{defn}
نظریهٔ~$\ZFminus$ این اصول موضوع را دارد: گسترش‌مندی، اجتماع، زوج‌ها،
مجموعه‌های توانی، نامتناهی و همهٔ نمونه‌های طرحواره‌های جداسازی و
جایگزینی. به بیان دیگر، $\ZFminus$ جایگزینی را به~$\Zminus$ می‌افزاید.
\end{defn}
\noindent
این نام نشان‌دهندهٔ نظریهٔ مجموعه‌های \emph{زرملو--فرنکل} است
(\emph{منهای} چیزی که بعداً به آن خواهیم رسید). افتخار نصیب فرنکل شده
است، زیرا صورت‌بندی جایگزینی در~\citeyear{Fraenkel1922} به او نسبت داده
می‌شود، هرچند نخستین صورت‌بندی دقیق را~\citet{Skolem1922} عرضه کرد.

\end{document}
